\ECchapter{05}{Energy Storage}{Why is energy storage essential to the energy transition?}{ECPurple}

\ECObjectives{
\begin{itemize}[leftmargin=*,itemsep=1pt,topsep=1pt]
\item Explain why electrical grids need storage and flexibility.
\item Distinguish power, stored energy and discharge duration.
\item Compare storage technologies using efficiency, response time, lifetime and scale.
\item Select a suitable storage mix for a given engineering problem.
\end{itemize}}

\begin{multicols}{2}

\begin{engineeringnews}
As solar and wind generation expands, grid operators are deploying batteries,
pumped-storage plants and other flexible resources. The objective is not to find one
universal storage device, but to combine technologies operating over seconds, hours, days
and seasons.
\end{engineeringnews}

\begin{ecarticle}
Electricity networks must keep production and consumption balanced at every instant. A
sudden mismatch changes grid frequency and may damage equipment or trigger automatic
disconnections. Storage systems absorb energy when production is greater than demand and
return it when the system requires additional power.

Engineers distinguish \textbf{power} from \textbf{energy}. Power, measured in watts,
describes how quickly a system can charge or discharge. Stored energy, measured in
watt-hours or joules, describes how long that power can be maintained. Their ratio gives an
approximate discharge duration:
\[
t=\frac{E}{P}.
\]
A 20~MW battery containing 40~MWh can therefore operate at full power for about two hours.

Lithium-ion batteries respond in fractions of a second and have a high round-trip
efficiency. They are suitable for frequency control, electric vehicles and shifting solar
energy from midday to the evening. Their limitations include ageing, fire protection, cost
and the supply of raw materials.

Pumped-storage hydroelectricity stores much larger amounts of energy. Water is pumped to an
upper reservoir and later released through a reversible turbine. Such plants can operate for
many decades, but they require suitable terrain and major civil-engineering works.

Compressed air, thermal storage and hydrogen can address longer durations. Hydrogen is
produced by electrolysis, stored, and later used in a fuel cell or turbine. Its round-trip
efficiency is lower because several energy conversions are required. It may nevertheless be
useful when storage duration and capacity matter more than efficiency, especially for
seasonal balancing.

No storage option is optimal for every task. A resilient low-carbon grid combines storage
with interconnections, demand response, controllable generation and accurate forecasting.
\end{ecarticle}

\begin{tcolorbox}[title=\sffamily\bfseries Did You Know?,colback=yellow!10,colframe=ECOrange]
A highly efficient storage system can still be unsuitable if its response time, discharge
duration, lifetime, safety requirements or location do not match the service required.
\end{tcolorbox}

\columnbreak

\begin{engineeringfocus}
\textbf{Pumped-storage hydroelectricity}

\begin{center}
\begin{tikzpicture}[scale=.68,transform shape,>=Latex,font=\scriptsize]
\fill[ECBlue!18] (0,2.7) rectangle (2.4,3.35);
\fill[ECBlue!18] (5.2,.1) rectangle (7.8,.75);
\draw[thick] (0,2.7) rectangle (2.4,3.35) node[midway]{upper reservoir};
\draw[thick] (5.2,.1) rectangle (7.8,.75) node[midway]{lower reservoir};
\node[draw=ECBlue,fill=ECLightBlue,rounded corners,minimum width=20mm,minimum height=9mm] (m) at (3.8,1.65) {motor--generator};
\node[draw=ECGreen,fill=ECGreen!10,rounded corners,minimum width=17mm,minimum height=9mm] (t) at (3.8,.65) {pump--turbine};
\draw[very thick,ECBlue] (2.4,3.0) -- (3.8,2.0) -- (t.north);
\draw[very thick,ECBlue] (t.east) -- (5.2,.45);
\draw[<->,very thick,ECOrange] (m) -- (t);
\draw[->,very thick,ECGreen] (2.8,2.35) -- node[above,sloped]{generation} (4.7,1.0);
\draw[->,very thick,ECPurple] (4.8,.82) -- node[below,sloped]{pumping} (2.7,2.2);
\end{tikzpicture}
\end{center}

During low demand, electrical energy drives the pump and raises water. During high demand,
the water flows downhill and drives the turbine. The stored gravitational energy is
$E=mgh$.
\end{engineeringfocus}

\begin{keyfigures}
\begin{tabularx}{\linewidth}{@{}Xr@{}}
Lithium-ion round-trip efficiency & 85--95\%\\
Pumped-storage hydro & 70--85\%\\
Hydrogen power-to-power & 30--45\%\\
Typical battery duration & 1--8 h\\
Seasonal storage & weeks--months
\end{tabularx}
\end{keyfigures}

\begin{tcolorbox}[title=\sffamily\bfseries Storage technologies compared,colback=white,colframe=ECBlue]
\centering
\begin{tikzpicture}
\begin{axis}[
width=\linewidth,height=47mm,
xbar,bar width=5pt,
xmin=0,xmax=100,
xlabel={Representative round-trip efficiency (\%)},
symbolic y coords={Hydrogen,Compressed air,Pumped hydro,Flywheel,Lithium-ion},
ytick=data,
yticklabel style={font=\scriptsize},
xticklabel style={font=\scriptsize},
label style={font=\scriptsize},
grid=major]
\addplot[fill=ECBlue] table[x=efficiency,y=technology,col sep=comma]
{data/EC05/storage_comparison.csv};
\end{axis}
\end{tikzpicture}

{\scriptsize Representative values only: actual performance depends on design and operating
conditions.}
\end{tcolorbox}

\begin{tcolorbox}[title=\sffamily\bfseries Engineering Insight,colback=white,colframe=ECOrange]
Storage sizing requires two independent decisions. The power rating must cover the largest
instantaneous imbalance, while the energy capacity must cover the required duration. A
system can have enough power but run empty too soon, or contain enough energy but deliver it
too slowly.
\end{tcolorbox}

\begin{vocabulary}
\begin{tabularx}{\linewidth}{@{}lX@{}}
round-trip efficiency & rendement aller-retour\\
charge / discharge & charge / décharge\\
storage duration & durée de stockage\\
reservoir & réservoir\\
flywheel & volant d'inertie\\
demand response & effacement de consommation\\
grid frequency & fréquence du réseau
\end{tabularx}
\end{vocabulary}

\begin{questions}
\item Why must electricity production and consumption remain balanced?
\item Distinguish power from stored energy.
\item A 12~MW system stores 48~MWh. What is its full-power discharge duration?
\item Why are batteries suitable for rapid grid services?
\item Explain the operating cycle of a pumped-storage plant.
\item Why can hydrogen be useful despite its lower efficiency?
\item Which criteria should engineers consider besides efficiency?
\end{questions}

\begin{challenge}
A remote island has a 12~MW peak demand, a 20~MW wind farm and a 10~MW solar farm.
Propose a storage strategy for frequency control, night-time supply and a five-day period of
weak wind. Select at least two technologies and justify their power, energy capacity,
duration, efficiency, safety and environmental impact. Present a block diagram and a table
of design choices.
\end{challenge}

\begin{applicationexercise}
A pumped-storage plant transfers $5.0\times10^6$~m$^3$ of water through a vertical height of
420~m. Estimate the stored gravitational energy using $E=\rho Vgh$, with
$\rho=1000~\mathrm{kg\,m^{-3}}$. Express the result in GWh, then estimate the recoverable
energy for a round-trip efficiency of 78\%. Finally, calculate how long the plant could
supply 1.0~GW at constant power.
\end{applicationexercise}

\begin{speaking}
In teams, defend one storage technology in a three-minute technical pitch. Explain where it
performs well, then identify one service for which it would be a poor choice.
\end{speaking}

\end{multicols}

\subsection*{Sources and further reading}
\ECsource{IEA -- Grid-scale storage}{https://www.iea.org/energy-system/electricity/grid-scale-storage}
\quad
\ECsource{NREL -- Energy storage}{https://www.nrel.gov/grid/energy-storage.html}
